%%%%%%%%%%%%%%%%%%%%%%%%%%%%%%%%%%%%%%%%%
% a0poster Landscape Poster
% LaTeX Template
% Deep Food Image Classifier with Evolutionary HPO
%
% Author: V. Harsha Vardhan Yellela
% Course: MCS-5993 Evolutionary Computing & Deep Learning
% Institution: Lawrence Technological University
%%%%%%%%%%%%%%%%%%%%%%%%%%%%%%%%%%%%%%%%%

\documentclass[a0,landscape]{a0poster}

\usepackage{listings}
\usepackage{color}
\usepackage{multicol}
\columnsep=50pt
\renewcommand{\columnseprulecolor}{\color{LightBlue}}
\columnseprule=3pt
\usepackage[svgnames]{xcolor}
\usepackage{times}
\usepackage{graphicx}
\graphicspath{{figures/}}
\usepackage{booktabs}
\usepackage[font=small,labelfont=bf]{caption}
\usepackage{amsfonts, amsmath, amsthm, amssymb}
\usepackage{wrapfig}
\usepackage{placeins}
\usepackage{titlesec}

% Reduce spacing throughout
\setlength{\parskip}{0pt}
\setlength{\baselineskip}{1.05\baselineskip}
\titlespacing*{\section}{0pt}{0.3em}{0.2em}

\begin{document}

%----------------------------------------------------------------------------------------
%	POSTER HEADER 
%----------------------------------------------------------------------------------------
\begin{minipage}[b]{0.55\linewidth}
\veryHuge \color{DarkOrange} \textbf{Deep Food Image Classifier: Comparing \\ Evolutionary Hyperparameter Optimization Methods} \color{Black}\\
\LARGE\\V Harsha Vardhan Yellela, MCS Candidate
\end{minipage}
%
\begin{minipage}[b]{0.25\linewidth}
\color{DarkSlateGray}\Large \textbf{Contact Information:}\\
Dept. of Mathematics and Computer Science\\
Lawrence Technological University\\
Science Building S121D\\
\end{minipage}
%
\begin{minipage}[b]{0.18\linewidth}
\includegraphics[width=20cm]{LTU2.png}
\end{minipage}

\vspace{1cm}

%----------------------------------------------------------------------------------------
\begin{multicols*}{4}

%----------------------------------------------------------------------------------------
%	ABSTRACT
%----------------------------------------------------------------------------------------
\color{Teal}
\section*{Abstract}
\color{black}
This work presents a comprehensive comparison of three hyperparameter optimization (HPO) methods for deep learning: hand-coded ES(1+1) Evolution Strategy, Keras Tuner Hyperband, and DEAP library-based ($\mu$+$\lambda$)-ES. We develop a 10-class American food image classifier using a custom ResNet-style CNN architecture with residual connections, achieving up to \textbf{79.8\% test accuracy}—a \textbf{24\% improvement} over the baseline. Our experiments demonstrate that evolutionary strategies can match or exceed gradient-free hyperparameter search methods while providing interpretable optimization trajectories. The system utilizes multi-GPU mixed-precision training on dual RTX 3090 GPUs, processes 10,000 images across 10 food categories, and optimizes 5 hyperparameters: learning rate, dropout rate, label smoothing, rotation range, and zoom range.

%----------------------------------------------------------------------------------------
%	INTRODUCTION
%----------------------------------------------------------------------------------------
\color{Teal}
\section*{Introduction}
\color{Black}
Hyperparameter optimization remains a critical challenge in deep learning, significantly impacting model performance. While grid search and random search are common approaches, evolutionary algorithms offer a principled way to explore the hyperparameter space efficiently.

\begin{center}\vspace{-0.5cm}
\includegraphics[width=0.85\linewidth]{food_samples.png}
\captionof{figure}{\color{Teal}Sample images from our 10-class American food dataset: chicken wings, churros, french fries, hamburger, hot dog, ice cream, mac \& cheese, pancakes, pizza, waffles}
\end{center}\vspace{-0.4cm}

\textbf{Research Questions:}
\begin{enumerate}
    \item Can hand-coded ES(1+1) match library-based HPO methods?
    \item How do different evolutionary strategies compare for CNN optimization?
    \item What hyperparameters have the most impact on food classification?
\end{enumerate}

Our food classification system processes images in real-time, supports out-of-scope detection via confidence thresholding, and achieves inference times of 3-5 seconds per image on consumer hardware.

%----------------------------------------------------------------------------------------
%	DATASET
%----------------------------------------------------------------------------------------
\color{Teal}
\section*{Dataset}
\color{Black}
\textbf{Source:} Food-101 dataset (subset of 10 American food classes)

\textbf{Statistics:} 10,000 images (1,000/class) $\bullet$ Train/Val/Test: 70\%/15\%/15\% $\bullet$ Size: 224$\times$224 RGB

\textbf{Food Classes:}
\begin{minipage}[t]{0.48\linewidth}
\begin{enumerate}
    \item Chicken Wings
    \item Churros
    \item French Fries
    \item Hamburger
    \item Hot Dog
\end{enumerate}
\end{minipage}
\hfill
\begin{minipage}[t]{0.48\linewidth}
\begin{enumerate}
\setcounter{enumi}{5}
    \item Ice Cream
    \item Macaroni \& Cheese
    \item Pancakes
    \item Pizza
    \item Waffles
\end{enumerate}
\end{minipage}

\textbf{Data Augmentation:} Horizontal flip $\bullet$ Random rotation $\bullet$ Random zoom $\bullet$ Contrast adjustment

%----------------------------------------------------------------------------------------
%	METHODOLOGY
%----------------------------------------------------------------------------------------
\color{Teal}
\section*{Methodology}
\color{Black}

\textbf{CNN Architecture:} Custom ResNet-style with 4 residual blocks

\begin{minipage}[t]{0.48\linewidth}
\begin{itemize}
    \item Initial Conv: 64 filters, 7$\times$7, stride 2
    \item ResBlock 1: 64 $\rightarrow$ 128 ch
    \item ResBlock 2: 128 $\rightarrow$ 256 ch
\end{itemize}
\end{minipage}
\hfill
\begin{minipage}[t]{0.48\linewidth}
\begin{itemize}
    \item ResBlock 3: 256 $\rightarrow$ 512 ch
    \item ResBlock 4: 512 $\rightarrow$ 1024 ch
    \item GAP + Dense(1024) + Softmax(10)
\end{itemize}
\end{minipage}

\textbf{Total Parameters:} $\sim$20.6M

\begin{center}\vspace{-0.3cm}
\includegraphics[width=0.85\linewidth]{cnn_architecture.png}
\captionof{figure}{\color{Teal}ResNet-style CNN Architecture with Residual Connections}
\end{center}\vspace{-0.3cm}

\textbf{Hyperparameters Optimized:}
\begin{center}
\begin{tabular}{|l|c|c|}
\hline
\textbf{Parameter} & \textbf{Range} & \textbf{Type} \\
\hline
Learning Rate & $[10^{-5}, 10^{-2}]$ & Log \\
Dropout Rate & $[0.1, 0.5]$ & Linear \\
Label Smoothing & $[0.0, 0.3]$ & Linear \\
Rotation/Zoom & $[0.0, 0.2/0.3]$ & Linear \\
\hline
\end{tabular}
\end{center}

%----------------------------------------------------------------------------------------
%	HPO METHODS
%----------------------------------------------------------------------------------------
\color{Teal}
\section*{HPO Methods Compared}
\color{Black}

\textbf{1. Baseline CNN:} Fixed HP (LR=0.002, Dropout=0.25), 50 epochs

\textbf{2. ES(1+1) - Hand-coded:}
\begin{itemize}
    \item Single parent, single offspring strategy
    \item Gaussian mutation with 1/5th success rule
    \item $\sigma$ adaptation: $\sigma \leftarrow \sigma \cdot e^{\pm 0.2}$
    \item 8 generations, 30 epochs per evaluation
\end{itemize}

\textbf{3. Keras Tuner (Hyperband):} Successive halving, Factor=3, 174 trials

\textbf{4. DEAP ($\mu$+$\lambda$)-ES:} Population $\mu$=5, $\lambda$=8, 5 generations

\begin{center}\vspace{-0.3cm}
\includegraphics[width=0.85\linewidth]{es_convergence.png}
\captionof{figure}{\color{Teal}ES(1+1) Fitness Convergence over Generations}
\end{center}\vspace{-0.3cm}

%----------------------------------------------------------------------------------------
%	RESULTS
%----------------------------------------------------------------------------------------
\color{Teal}
\section*{Results}
\color{Black}

\textbf{Test Accuracy Comparison:}
\begin{center}
\begin{tabular}{|l|c|c|}
\hline
\textbf{Method} & \textbf{Test Acc.} & \textbf{Improvement} \\
\hline
Baseline CNN & 64.37\% & --- \\
\hline
ES(1+1) & \textbf{79.80\%} & \textbf{+23.97\%} \\
\hline
Keras Tuner & 79.27\% & +23.14\% \\
\hline
DEAP ($\mu$+$\lambda$) & 79.33\% & +23.25\% \\
\hline
\end{tabular}
\captionof{table}{\color{Teal}Performance Comparison of All Methods}
\end{center}

\begin{center}\vspace{-0.3cm}
\includegraphics[width=0.85\linewidth]{accuracy_comparison.png}
\captionof{figure}{\color{Teal}Test Accuracy Comparison: All Four Methods}
\end{center}\vspace{-0.3cm}

\textbf{Optimized Hyperparameters:}
\begin{center}
\begin{tabular}{|l|c|c|c|c|}
\hline
\textbf{HP} & \textbf{Base} & \textbf{ES(1+1)} & \textbf{Tuner} & \textbf{DEAP} \\
\hline
LR & 2e-3 & 1.38e-3 & 1.41e-3 & 9.86e-5 \\
\hline
Dropout & 0.25 & 0.256 & 0.30 & 0.10 \\
\hline
Label Sm. & 0.05 & 0.30 & 0.28 & 0.077 \\
\hline
Rotation & 0.05 & 0.066 & 0.05 & 0.031 \\
\hline
Zoom & 0.10 & 0.114 & 0.10 & 0.327 \\
\hline
\end{tabular}
\captionof{table}{\color{Teal}Optimized Hyperparameters by Method}
\end{center}

\begin{center}\vspace{-0.3cm}
\includegraphics[width=0.85\linewidth]{hyperparameter_comparison.png}
\captionof{figure}{\color{Teal}Hyperparameter Values Across Methods}
\end{center}\vspace{-0.3cm}

\textbf{Key Findings:}
\begin{itemize}
    \item All three HPO methods achieved similar performance ($\sim$79\%)
    \item Hand-coded ES(1+1) slightly outperformed library methods
    \item Label smoothing had significant impact (0.05 $\rightarrow$ 0.28-0.30)
    \item Learning rate convergence: $\sim$1.4$\times$10$^{-3}$ optimal
\end{itemize}

%----------------------------------------------------------------------------------------
%	TRAINING DETAILS
%----------------------------------------------------------------------------------------
\color{Teal}
\section*{Training Configuration}
\color{Black}

\begin{minipage}[t]{0.48\linewidth}
\textbf{Hardware:}
\begin{itemize}
    \item 2$\times$ NVIDIA RTX 3090 (48GB)
    \item 21GB System RAM + 8GB Swap
    \item WSL2 Ubuntu environment
\end{itemize}
\end{minipage}
\hfill
\begin{minipage}[t]{0.48\linewidth}
\textbf{Software:}
\begin{itemize}
    \item TensorFlow 2.10 with Keras
    \item Mixed Precision (float16)
    \item MirroredStrategy for multi-GPU
\end{itemize}
\end{minipage}

\textbf{Training Time:} Baseline: 1.5h $\bullet$ ES(1+1): 6-8h (8 gen) $\bullet$ Keras Tuner: 6h (174 trials) $\bullet$ DEAP: 4h (5 gen)

\begin{center}\vspace{-0.3cm}
\includegraphics[width=0.85\linewidth]{training_curves.png}
\captionof{figure}{\color{Teal}Training and Validation Accuracy Curves}
\end{center}\vspace{-0.3cm}

%----------------------------------------------------------------------------------------
%	INFERENCE
%----------------------------------------------------------------------------------------
\color{Teal}
\section*{Inference \& Deployment}
\color{Black}

\textbf{Out-of-Scope Detection:} Confidence threshold: 60\% $\bullet$ Predictions below threshold flagged as ``out of scope'' $\bullet$ Prevents misclassification of non-food images

\textbf{Model Deployment:}
\begin{itemize}
    \item Hugging Face Hub: \texttt{HAR5HA-YELLELA/american\_food\_classifier}
    \item Inference: 3-5 seconds $\bullet$ Memory: 4-5 GB VRAM
\end{itemize}

\begin{center}\vspace{-0.3cm}
\includegraphics[width=0.85\linewidth]{inference_example.png}
\captionof{figure}{\color{Teal}Inference Example: Churros classified with 75.1\% confidence}
\end{center}\vspace{-0.3cm}

%----------------------------------------------------------------------------------------
%	CONCLUSION
%----------------------------------------------------------------------------------------
\color{Teal}
\section*{Conclusion}
\color{Black}

\begin{itemize}
    \item \textbf{All HPO methods effective:} 23-24\% improvement over baseline
    \item \textbf{ES(1+1) competitive:} Hand-coded implementation matched library methods
    \item \textbf{Label smoothing critical:} Increased from 0.05 to 0.28-0.30
    \item \textbf{Evolutionary strategies viable:} For CNN hyperparameter optimization
    \item \textbf{Practical deployment:} Model available on Hugging Face Hub
\end{itemize}

\textbf{Recommendations:} 1. Use ES(1+1) for simplicity 2. Keras Tuner for extensive search 3. DEAP for complex multi-objective optimization

%----------------------------------------------------------------------------------------
%	FUTURE WORK
%----------------------------------------------------------------------------------------
\color{Teal}
\section*{Future Work}
\color{Black}

\begin{enumerate}
    \item Extend to full Food-101 dataset (101 classes)
    \item Implement CMA-ES for covariance matrix adaptation
    \item Multi-objective optimization (accuracy vs. inference time)
    \item Neural Architecture Search (NAS) with evolutionary methods
\end{enumerate}

%----------------------------------------------------------------------------------------
%	REFERENCES
%----------------------------------------------------------------------------------------
\color{Teal}
\section*{References}
\color{Black}
\small
\begin{enumerate}
    \item Bossek, J. et al. (2019). "A Survey on Hyperparameter Optimization." \textit{arXiv:1907.10902}
    \item Li, L. et al. (2018). "Hyperband: A Novel Bandit-Based Approach." \textit{JMLR}
    \item Fortin, F.A. et al. (2012). "DEAP: Evolutionary Algorithms Made Easy." \textit{JMLR}
    \item He, K. et al. (2016). "Deep Residual Learning for Image Recognition." \textit{CVPR}
\end{enumerate}

\color{Teal}
\end{multicols*}
\end{document}
